\documentclass[11pt,a4paper]{article}
\usepackage[utf8]{inputenc}
\usepackage[T1]{fontenc}
\usepackage[margin=2.5cm]{geometry}
\usepackage{amsmath,amssymb,bm}
\usepackage{graphicx}

\setlength{\parindent}{0pt}
\setlength{\parskip}{0.6em}

\title{Practical Work 2 -- Image Deblurring using Gibbs Sampling\\Answers}
\author{}
\date{}

\newcommand{\Dt}{\tilde{\mathbf D}}

\begin{document}
\maketitle

\section*{Question 1}
\begin{itemize}
  \item $\tau_e$: noise precision (inverse of the noise variance). A large $\tau_e$ means a weakly noisy observation.
  \item $\tau_x$: precision of the prior (regularization weight). A large $\tau_x$ enforces a strongly smooth image.
  \item $\mathbf D$: finite-difference operator (discrete gradient). $\|\mathbf D\mathbf x\|^2$ sums the squared differences between neighbouring pixels, so the prior penalizes abrupt variations and favours smooth images.
\end{itemize}

\section*{Question 2}
The acquired image $\mathbf y$ is blurred: edges are smeared and fine details and textures are lost compared with the true image $\mathbf x$. It is also corrupted by additive noise, which is of low level here.

\section*{Question 3}
The kernel is the point spread function (PSF): the response of the system to a point source. It is a $7\times 7$, symmetric, non-negative kernel summing to $1$, peaked at its centre. It defines $\mathbf H$: the acquired image is the convolution of the true image with it. It acts as a low-pass filter, each pixel being replaced by a weighted average of its neighbours.

\section*{Question 4}
c.f.\ code (\texttt{Fourier\_D.m}):
\[
\Dt = |\widehat{\mathbf D}_x|^2 + |\widehat{\mathbf D}_y|^2 .
\]

Usefulness: $\mathbf D^\top\mathbf D$ (like $\mathbf H^\top\mathbf H$) is a convolution operator, hence diagonalized by the DFT. The precision matrix of $\mathbf x \mid \tau_e,\tau_x,\mathbf y$,
\[
\tau_e\mathbf H^\top\mathbf H + \tau_x\mathbf D^\top\mathbf D \;\longrightarrow\; \tau_e|\widehat{\mathbf H}|^2 + \tau_x\Dt ,
\]
becomes an element-wise expression. Its inversion and the sampling of $\mathbf x$ then cost $\mathcal O(N\log N)$ instead of inverting an $N\times N$ matrix. It is also used to compute $\|\mathbf D\mathbf x\|^2$ in the $\tau_x$ update (Parseval).

\section*{Derivation 4.1 -- Conditional distribution of $\tau_e$}
\textbf{Goal:} keep only the terms depending on $\tau_e$ and recognize the kernel of a Gamma distribution, then identify $\alpha_e$ and $\beta_e$.

\textbf{1. Posterior} (Bayes' rule):
\[
p(\mathbf x,\tau_e,\tau_x\mid\mathbf y) \propto
\underbrace{p(\mathbf y\mid\mathbf x,\tau_e)}_{\text{Gaussian likelihood}}\;
\underbrace{p(\mathbf x\mid\tau_x)}_{\text{Gaussian prior}}\;
\underbrace{p(\tau_e)\,p(\tau_x)}_{\text{Gamma priors}} .
\]

\textbf{2. Conditional:} $p(\tau_e\mid\mathbf x,\tau_x,\mathbf y)\propto p(\mathbf x,\tau_e,\tau_x\mid\mathbf y)$. The factors $p(\mathbf x\mid\tau_x)$ and $p(\tau_x)$ do not depend on $\tau_e$, so
\[
p(\tau_e\mid\mathbf x,\mathbf y) \propto p(\mathbf y\mid\mathbf x,\tau_e)\,p(\tau_e).
\]

\textbf{3. Likelihood:} $\mathbf y\mid\mathbf x,\tau_e\sim\mathcal N(\mathbf H\mathbf x,\tau_e^{-1}\mathbf I_N)$, hence
\[
p(\mathbf y\mid\mathbf x,\tau_e) = \left(\frac{\tau_e}{2\pi}\right)^{N/2}\exp\!\left(-\frac{\tau_e}{2}\|\mathbf y-\mathbf H\mathbf x\|^2\right)
\propto \tau_e^{N/2}\exp\!\left(-\frac{\tau_e}{2}\|\mathbf y-\mathbf H\mathbf x\|^2\right).
\]

\textbf{4. Prior} (shape--rate): $\tau_e\sim\Gamma(\alpha_{e0},\beta_{e0})$, i.e.\ $p(\tau_e)\propto\tau_e^{\alpha_{e0}-1}\exp(-\beta_{e0}\tau_e)$.

\textbf{5. Product:}
\[
p(\tau_e\mid\mathbf x,\mathbf y) \propto
\tau_e^{\,\alpha_{e0}+\frac N2-1}\exp\!\left[-\tau_e\left(\beta_{e0}+\frac12\|\mathbf y-\mathbf H\mathbf x\|^2\right)\right].
\]

\textbf{6. Identification:} this is the kernel of $\Gamma(\alpha,\beta)\propto\tau^{\alpha-1}e^{-\beta\tau}$, hence
\[
\boxed{\;p(\tau_e\mid\mathbf x,\mathbf y)=\Gamma(\tau_e;\alpha_e,\beta_e),\qquad
\alpha_e=\alpha_{e0}+\frac N2,\qquad
\beta_e=\beta_{e0}+\frac12\|\mathbf y-\mathbf H\mathbf x\|^2\;}
\]

\textbf{Remark (implementation):} with the unnormalized DFT (\texttt{fft2}), Parseval gives
$\|\mathbf y-\mathbf H\mathbf x\|^2=\frac1N\sum|\widehat{\mathbf y}-\widehat{\mathbf H}\odot\widehat{\mathbf x}|^2$.

\section*{Question 5}
c.f.\ code (\texttt{main.m}, block (a)):
\begin{verbatim}
alpha_e = alpha_e0 + Npix/2;
res2    = sum(sum(abs(y_hat - Hhat .* x_hat).^2)) / Npix;  % ||y-Hx||^2
beta_e  = beta_e0 + res2/2;
tau_e   = rndgamma_shape_rate(alpha_e, beta_e);
\end{verbatim}

\section*{Derivation 4.2 -- Conditional distribution of $\tau_x$}
\textbf{1. Conditional:} the likelihood $p(\mathbf y\mid\mathbf x,\tau_e)$ and $p(\tau_e)$ do not depend on $\tau_x$, so
\[
p(\tau_x\mid\mathbf x,\mathbf y) \propto p(\mathbf x\mid\tau_x)\,p(\tau_x).
\]

\textbf{2. Prior on $\mathbf x$:} Gaussian with precision matrix $\tau_x\mathbf D^\top\mathbf D$. Its normalization constant depends on $\tau_x$ and must be kept:
\[
p(\mathbf x\mid\tau_x) \propto \tau_x^{N/2}\exp\!\left(-\frac{\tau_x}{2}\|\mathbf D\mathbf x\|^2\right).
\]

\textbf{3. Hyperprior:} $p(\tau_x)\propto\tau_x^{\alpha_{x0}-1}\exp(-\beta_{x0}\tau_x)$.

\textbf{4. Product:}
\[
p(\tau_x\mid\mathbf x,\mathbf y) \propto
\tau_x^{\,\alpha_{x0}+\frac N2-1}\exp\!\left[-\tau_x\left(\beta_{x0}+\frac12\|\mathbf D\mathbf x\|^2\right)\right],
\]
which is the kernel of a Gamma distribution: $\tau_x\mid\mathbf x,\mathbf y\sim\Gamma(\alpha_x,\beta_x)$.

\section*{Question 6}
Identifying the exponent $\alpha_x-1=\alpha_{x0}+\frac N2-1$ and the rate:
\[
\boxed{\;\alpha_x=\alpha_{x0}+\frac N2,\qquad \beta_x=\beta_{x0}+\frac12\|\mathbf D\mathbf x\|^2\;}
\]
In the Fourier domain (Parseval): $\|\mathbf D\mathbf x\|^2=\frac1N\sum\Dt\odot|\widehat{\mathbf x}|^2$.

c.f.\ code (\texttt{main.m}, block (b)):
\begin{verbatim}
xnormD2 = sum(sum(D .* abs(x_hat).^2)) / Npix;  % ||Dx||^2
alpha_x = alpha_x0 + Npix/2;
beta_x  = beta_x0 + xnormD2/2;
tau_x   = rndgamma_shape_rate(alpha_x, beta_x);
\end{verbatim}

\section*{Question 8}
With periodic boundary conditions, $\mathbf H^\top\mathbf H$ and $\mathbf D^\top\mathbf D$ are circulant (convolutions), hence diagonalized by the DFT. The posterior precision $\tau_e\mathbf H^\top\mathbf H+\tau_x\mathbf D^\top\mathbf D$ becomes the diagonal $\tau_e|\widehat{\mathbf H}|^2+\tau_x\Dt$:
\begin{itemize}
  \item inversion and sampling are element-wise, in $\mathcal O(N)$, plus FFTs in $\mathcal O(N\log N)$;
  \item in the spatial domain, they would require inverting or factorizing (Cholesky) an $N\times N$ matrix, in $\mathcal O(N^3)$. Here $N=256^2=65\,536$, i.e.\ about $4.3\cdot10^9$ entries ($\approx 34$\,GB in double precision).
\end{itemize}

\section*{Question 9}
c.f.\ code (\texttt{main.m}, block (c)):
\begin{verbatim}
num   = tau_e * conj(Hhat) .* y_hat;
denom = tau_e * abs(Hhat).^2 + tau_x * D;
muhat = num ./ denom;
\end{verbatim}
$\widehat{\mathbf H}^*$ is the Fourier counterpart of $\mathbf H^\top$. Here the PSF is symmetric, so $\widehat{\mathbf H}$ is real and $\widehat{\mathbf H}^*=\widehat{\mathbf H}$. At frequency $0$, $\widehat{\mathbf H}=1$ and $\Dt=0$, so $\text{denom}=\tau_e>0$: there is no division by zero.

\section*{Question 10}
c.f.\ code (\texttt{main.m}, block (c)):
\begin{verbatim}
U = fft2(randn(M,N))/sqrt(Npix);
x_hat_s = muhat + sqrt(Npix) * U ./ sqrt(denom);
x = real(ifft2(x_hat_s));
\end{verbatim}
With $\mathbf W=\mathbf F/\sqrt N$ the unitary DFT and $\mathbf z\sim\mathcal N(\mathbf 0,\mathbf I)$, $\mathbf x=\boldsymbol\mu+\mathbf W^H\widehat{\boldsymbol\Sigma}^{-1/2}\mathbf W\mathbf z$ has covariance $\mathbf W^H\widehat{\boldsymbol\Sigma}^{-1}\mathbf W=(\tau_e\mathbf H^\top\mathbf H+\tau_x\mathbf D^\top\mathbf D)^{-1}$. With the unnormalized \texttt{fft2}, this gives $\widehat{\mathbf x}=\widehat{\boldsymbol\mu}+\mathbf F\mathbf z\oslash\sqrt{\widehat{\boldsymbol\Sigma}}=\widehat{\boldsymbol\mu}+\sqrt N\,\mathbf U\oslash\sqrt{\widehat{\boldsymbol\Sigma}}$, hence the $\sqrt N$ factor. This was checked by Monte-Carlo on the variance of one pixel. The sample is not named \texttt{X}, because \texttt{X} already holds the true image used for the PSNR.

\section*{Question 12}
The Burnin Method is used to discard the initial samples of the Markov chain, which may not represent the target distribution well.
However, on this image at 1k iterations keeping the first iterations does not change the results significantly. 28.63 db (with burnin) vs 28.62 db (without burnin).


\section*{Question 13}
Around 10 iterations are enough to reach a stable PSNR.
\begin{figure}[h!]
\centering
\includegraphics[width=\linewidth]{chains.png}
\caption{Evolution of $\tau_e$, $\tau_x$ and $\sigma_e=1/\sqrt{\tau_e}$ over the iterations (red dashed line: burn-in).}
\label{fig:chains}
\end{figure}

\section*{Question 14}
c.f code (\texttt{main.m}, block (d)):
If the iteration number is after the burn in AND if the index of the iteration is a multiple of the thinning factor (even because thin ==2), then the sample is kept. Otherwise it is discarded.


The condition \texttt{if it > BurnIn \&\& mod(it - BurnIn, Thin) == 0} allows selecting and storing the image samples only after the burn-in period, while applying a thinning step to limit their statistical correlation. Then, the code block computes the posterior mean estimator $\hat{x}$ (represented as \texttt{Xbar}) by calculating the empirical mean of all the preserved samples:   
\[
\hat{x} = \frac{1}{K} \sum_{k=1}^{K} x^{(k)}
\]


\section*{Question 15}

The reconstruccted image is granier than the original image. The grayscale is smoother. The black is lighter ans White is darker. 

\section*{Question 16}

It adapt with the input image. It automaticately converge to the best solution.
If we choose it manually we have to try different values and see which one is the best. It is not automatic and it is not optimal.

\section*{Question 17}

\begin{itemize}
  \item \textbf{Deconvolution}. We reverse the effect of the blurring convolution.
  \item \textbf{Bayesian}. x, $\tau_e$ and $\tau_x$ are random variables with prior laws. We compute this law and we sample it with Gibbs.
  \item \textbf{Unsupervised}. The hyperparameters $\tau_e$ and $\tau_x$ are estimated from the data, instead of being fixed manually.
\end{itemize}
\end{document}
